\chapter{总结与展望}\label{chap:conclusion}

\section{研究工作总结}\label{sec:summary-work}

围绕GBM候选药物分子的计算生成与优化问题，本文构建了GBM-PharmaMind候选药物分子生成框架。该框架以多模态大语言模型和分子图—文本生成模型为基础，将GBM领域知识输入、CoT-ToT协同推理和物理评估反馈统一到同一生成流程中。本文工作可从知识层面、框架层面和评估层面加以概括。

\subsection{GBM多模态知识库构建与注入}\label{sec:work-km}

在候选分子生成任务中，模型能否获得足够明确的疾病背景信息，会直接影响生成结果的针对性。本文围绕GBM药物设计场景，整理了与候选分子生成相关的靶点信息、临床治疗经验、参考分子结构以及CNS药物设计规则，使生成过程能够同时参考GBM疾病背景、血脑屏障相关约束和已知结构经验，从而减少完全依赖通用化学空间采样带来的不确定性。

具体来看，本文构建的知识库并不是单独作为背景资料保存，而是参与到Prompt构造和后续生成控制过程中。GBM靶点信息库用于提供EGFR、VEGFR2、IDH1等靶点的基本背景和设计提示；临床治疗数据库用于补充标准治疗方案、耐药机制和已有治疗经验；分子参考数据库用于提供已知药物或候选分子的结构参照；结构知识库则用于记录BBB穿透相关策略、PAINS警告结构以及基础类药性规则。分散的疾病知识和药物化学经验被转化为生成模型可以读取和利用的约束信息，为后续候选分子生成提供了较稳定的输入基础。

\subsection{CoT-ToT协同推理与分子生成模块}\label{sec:work-cot-tot}

本文在分子生成阶段引入CoT-ToT协同推理，主要是为了避免模型一次性输出完整SMILES时缺少中间控制的问题。GBM候选分子生成需要同时考虑靶点背景、BBB穿透倾向、结构有效性和成药性约束，若直接生成最终分子，模型的骨架选择、片段修饰和结构表达会被混合在同一输出中，后续难以判断问题出现在设计思路还是结构落地环节。因此，本文将生成过程拆分为骨架提议、片段组装和SMILES输出三个阶段，使候选分子的形成过程更加清晰。

在这一过程中，CoT用于组织分子设计的推理顺序，使模型先分析靶点和约束条件，再逐步形成候选结构；ToT用于在不同阶段保留多个候选方向，使生成过程不局限于单一路径。骨架提议阶段关注分子的主体结构和潜在药效团方向，片段组装阶段在已有骨架上进行局部修饰，SMILES输出阶段则将前两步得到的结构设想转化为可计算的分子表示。物理评估模块的结果会被写回后续Prompt，使模型能够根据上一轮生成结果调整后续候选方向。通过这种方式，本文将结构化推理、多分支搜索和外部评价反馈连接到同一生成流程中，提高了候选分子生成过程的可控性。

\subsection{基于结构警报检测与Clark公式的成药性快速评估体系}\label{sec:work-physics}

评估层面，本文建立了结合结构警报检测、BBB经验预测、Vina对接、MPO Reward和TotalScore的计算评价体系。生成阶段的评价结果用于剪枝与反馈，最终阶段则通过TotalScore对候选分子进行综合排序。该体系将药物化学规则和外部计算工具转化为可执行的筛选流程，并通过硬剪枝与Delta Scoring机制在安全性约束和化学空间探索之间进行折中。

\section{主要结论}\label{sec:conclusions}

在EGFR、VEGFR2和IDH1等代表性GBM靶点上的实验结果表明，GBM-PharmaMind能够生成具有一定预测活性、BBB穿透倾向和成药性特征的候选分子。Full Model在DrugCLIP、Vina对接分数和BBB相关指标上整体优于消融模型，说明领域知识输入、LoRA适配、CoT-ToT推理和物理评估反馈均对最终生成质量具有显著贡献。

\section{存在的不足}\label{sec:limitations}

本文主要从计算方法角度构建GBM候选分子生成与优化框架，整体流程已经覆盖知识输入、分阶段生成、物理评估反馈和最终候选排序等环节，但仍存在一定改进空间。首先，当前实验主要围绕EGFR、VEGFR2和IDH1等代表性靶点展开，尚未覆盖更多GBM相关分子亚型和耐药机制，因此框架在更复杂靶点组合下的稳定性仍需进一步分析。其次，最终评估中的Vina Score、ADMET性质等指标主要DrugClip和SwissADME等外部工具，对不同预测模型之间的一致性和适用边界讨论仍不充分。最后，ToT搜索宽度和深度目前采用固定设置，尚未根据中间评价结果动态调整，这在一定程度上限制了搜索过程的自适应能力。

\section{未来研究展望}\label{sec:future}

围绕本文工作，后续研究可以继续从框架扩展、评价体系和搜索策略三个方向推进。在框架扩展方面，可将更多GBM相关通路、分子亚型和突变背景纳入实验，并扩大候选分子生成规模，以检验GBM-PharmaMind在不同任务条件下的稳定性。在评价体系方面，可进一步引入更精细的BBB预测模型、P-gp风险评估和逆合成可行性分析，减少单一经验指标对最终排序结果的影响。在搜索策略方面，后续可探索动态ToT宽度和深度设置，使生成过程能够根据中间评分结果自适应调整搜索空间，从而提高候选分子生成的效率与质量。

此外，GBM-PharmaMind框架也可以进一步扩展到其他中枢神经系统疾病或多目标药物设计任务中。由于该框架将疾病知识、分子生成和外部评价反馈进行了模块化组织，未来只需替换相应的疾病知识库和评价约束，就可以适配新的候选分子生成场景。这种扩展方式有助于提升框架的通用性，也为多模态大语言模型在计算药物设计中的应用提供了进一步研究空间。
